% GENERATED FILE. Edit canonical lessons, then run npm run generate:books.
% canonical-source-hash: fnv1a64:d43e9761b7300bd0
% canonical-lessons: LA-C386-societas, LA-C386-humanitas, LA-C386-mos, LA-C386-institutum, LA-C386-concordia

\chapter{Living Together}
\label{ch:la-living-together}
\hlchaptermodality{386}

\begin{chapteropening}
\textbf{By the end of this chapter:} \emph{I can talk about a partnership, kindness and culture, custom and character, an institution and harmony.}

Complete the last lesson of chapter 386: I can talk about a partnership, kindness and culture, custom and character, an institution and harmony.
\end{chapteropening}

\section[societās, societātis]{societās, societātis — a partnership, an alliance; society}

\label{lesson:LA-C386-societas}



\begin{quote}
Before the new one: say the Latin for command, then the Latin for the first.
\end{quote}

\subsection*{You'll want to know: societās, societātis}
\textbf{societās, societātis} — \textquotedblleft{}a partnership, an alliance; society\textquotedblright{}.

\begin{quote}
Cum frātre societātem fēcit.
\end{quote}

\begin{quote}
\emph{He went into partnership with his brother.}
\end{quote}

\emph{societās, societātis} f., built on \emph{socius}, a partner, an ally. A business partnership in Roman law. English \emph{society}.

\begin{cousinweb}
\emph{societās} is built on \emph{socius} \textquotedblleft{}a partner, an ally\textquotedblright{} (EtymDB; Lewis \& Short), a word kin to \emph{sequor} \textquotedblleft{}to follow\textquotedblright{}: one who goes along with you. English \emph{society} comes from it.
\end{cousinweb}

\begin{grammarlens}[title={What you've built}]
One more word for telling, explaining and describing at B1.
\end{grammarlens}

\subsection*{Your turn}
\begin{itemize}
\raggedright
  \item \emph{Say it:} \emph{societās, societātis}
  \item \emph{Say it:} \emph{societās, societātis}, once more
  \item \emph{Say it:} a sentence of your own with \emph{societās, societātis}
  \item \emph{Recall:} say the Latin for command, then the Latin for the first, then say \emph{societās, societātis} again
\end{itemize}

\begin{tcolorbox}[breakable,colback=teal!4,colframe=teal!35!black,title={Before you move on}]
What does societās, societātis mean? (\textquotedblleft{}A partnership, an alliance; society\textquotedblright{}.) Say it once more.
\end{tcolorbox}
\section[hūmānitās, hūmānitātis]{hūmānitās, hūmānitātis — human nature; kindness; culture}

\label{lesson:LA-C386-humanitas}



\begin{quote}
Before the new one: say the Latin for the first, then the Latin for a partnership.
\end{quote}

\subsection*{You'll want to know: hūmānitās, hūmānitātis}
\textbf{hūmānitās, hūmānitātis} — \textquotedblleft{}human nature; kindness; culture\textquotedblright{}.

\begin{quote}
Hospitem summā hūmānitāte accēpit.
\end{quote}

\begin{quote}
\emph{He received the guest with the greatest kindness.}
\end{quote}

\emph{hūmānitās, hūmānitātis} f., built on \emph{hūmānus}, human. The kindness and learning that make a person truly human. English \emph{humanity}.

\begin{cousinweb}
\emph{hūmānitās} is built on \emph{hūmānus} \textquotedblleft{}human\textquotedblright{} (Lewis \& Short), an adjective kin to \emph{homō} \textquotedblleft{}a human being\textquotedblright{}. English \emph{humanity} comes from it.
\end{cousinweb}

\begin{grammarlens}[title={What you've built}]
One more word for telling, explaining and describing at B1.
\end{grammarlens}

\subsection*{Your turn}
\begin{itemize}
\raggedright
  \item \emph{Say it:} \emph{hūmānitās, hūmānitātis}
  \item \emph{Say it:} \emph{hūmānitās, hūmānitātis}, once more
  \item \emph{Say it:} a sentence of your own with \emph{hūmānitās, hūmānitātis}
  \item \emph{Recall:} say the Latin for the first, then the Latin for a partnership, then say \emph{hūmānitās, hūmānitātis} again
\end{itemize}

\begin{tcolorbox}[breakable,colback=teal!4,colframe=teal!35!black,title={Before you move on}]
What does hūmānitās, hūmānitātis mean? (\textquotedblleft{}Human nature; kindness; culture\textquotedblright{}.) Say it once more.
\end{tcolorbox}
\section[mōs, mōris]{mōs, mōris — a custom, a way; (plural) character, morals}

\label{lesson:LA-C386-mos}



\begin{quote}
Before the new one: say the Latin for a partnership, then the Latin for human nature.
\end{quote}

\subsection*{You'll want to know: mōs, mōris}
\textbf{mōs, mōris} — \textquotedblleft{}a custom, a way; (plural) character, morals\textquotedblright{}.

\begin{quote}
Hic mōs antīquus est.
\end{quote}

\begin{quote}
\emph{This is an old custom.}
\end{quote}

\emph{mōs, mōris} m. The plural \emph{mōrēs} is a person's ways, his character. \emph{Mōs maiōrum}, the custom of the ancestors. English \emph{moral} and \emph{morals}.

\begin{cousinweb}
The origin of \emph{mōs} is not settled; Lewis \& Short give no sure account. English \emph{moral} comes from its adjective \emph{mōrālis}.
\end{cousinweb}

\begin{grammarlens}[title={What you've built}]
One more word for telling, explaining and describing at B1.
\end{grammarlens}

\subsection*{Your turn}
\begin{itemize}
\raggedright
  \item \emph{Say it:} \emph{mōs, mōris}
  \item \emph{Say it:} \emph{mōs, mōris}, once more
  \item \emph{Say it:} a sentence of your own with \emph{mōs, mōris}
  \item \emph{Recall:} say the Latin for a partnership, then the Latin for human nature, then say \emph{mōs, mōris} again
\end{itemize}

\begin{tcolorbox}[breakable,colback=teal!4,colframe=teal!35!black,title={Before you move on}]
What does mōs, mōris mean? (\textquotedblleft{}A custom, a way; (plural) character, morals\textquotedblright{}.) Say it once more.
\end{tcolorbox}
\section[īnstitūtum, īnstitūtī]{īnstitūtum, īnstitūtī — an institution, a custom; a purpose}

\label{lesson:LA-C386-institutum}



\begin{quote}
Before the new one: say the Latin for human nature, then the Latin for a custom.
\end{quote}

\subsection*{You'll want to know: īnstitūtum, īnstitūtī}
\textbf{īnstitūtum, īnstitūtī} — \textquotedblleft{}an institution, a custom; a purpose\textquotedblright{}.

\begin{quote}
Īnstitūta patrum servāmus.
\end{quote}

\begin{quote}
\emph{We keep the institutions of our fathers.}
\end{quote}

\emph{īnstitūtum, īnstitūtī} n.: a thing set up and kept, a custom. English \emph{institute} and \emph{institution}.

\begin{cousinweb}
\emph{īnstitūtum} is the participle of \emph{īnstituō} \textquotedblleft{}to set up\textquotedblright{} (Lewis \& Short), from \emph{in} and \emph{statuō} \textquotedblleft{}to set, to place\textquotedblright{}. English \emph{institute} comes from it.
\end{cousinweb}

\begin{grammarlens}[title={What you've built}]
One more word for telling, explaining and describing at B1.
\end{grammarlens}

\subsection*{Your turn}
\begin{itemize}
\raggedright
  \item \emph{Say it:} \emph{īnstitūtum, īnstitūtī}
  \item \emph{Say it:} \emph{īnstitūtum, īnstitūtī}, once more
  \item \emph{Say it:} a sentence of your own with \emph{īnstitūtum, īnstitūtī}
  \item \emph{Recall:} say the Latin for human nature, then the Latin for a custom, then say \emph{īnstitūtum, īnstitūtī} again
\end{itemize}

\begin{tcolorbox}[breakable,colback=teal!4,colframe=teal!35!black,title={Before you move on}]
What does īnstitūtum, īnstitūtī mean? (\textquotedblleft{}An institution, a custom; a purpose\textquotedblright{}.) Say it once more.
\end{tcolorbox}
\section[concordia, concordiae]{concordia, concordiae — harmony, concord}

\label{lesson:LA-C386-concordia}



\begin{quote}
Before the new one: say the Latin for a custom, then the Latin for an institution.
\end{quote}

\subsection*{You'll want to know: concordia, concordiae}
\textbf{concordia, concordiae} — \textquotedblleft{}harmony, concord\textquotedblright{}.

\begin{quote}
Concordiā parvae rēs crēscunt.
\end{quote}

\begin{quote}
\emph{Through harmony small states grow.}
\end{quote}

\emph{concordia, concordiae} f., built on \emph{cor, cordis}, heart: hearts together. English \emph{concord}.

\begin{cousinweb}
\emph{concordia} is built on \emph{concors} \textquotedblleft{}of one heart\textquotedblright{} (Lewis \& Short), from \emph{con-} \textquotedblleft{}together\textquotedblright{} and \emph{cor, cordis} \textquotedblleft{}heart\textquotedblright{}. English \emph{concord} comes from it.
\end{cousinweb}

\begin{grammarlens}[title={What you've built}]
One more word for telling, explaining and describing at B1.
\end{grammarlens}

\subsection*{Your turn}
\begin{itemize}
\raggedright
  \item \emph{Say it:} \emph{concordia, concordiae}
  \item \emph{Say it:} \emph{concordia, concordiae}, once more
  \item \emph{Say it:} a sentence of your own with \emph{concordia, concordiae}
  \item \emph{Recall:} say the Latin for a custom, then the Latin for an institution, then say \emph{concordia, concordiae} again
\end{itemize}

\begin{tcolorbox}[breakable,colback=teal!4,colframe=teal!35!black,title={Before you move on}]
What does concordia, concordiae mean? (\textquotedblleft{}Harmony, concord\textquotedblright{}.) Say it once more.
\end{tcolorbox}
